% ===========================================================================
% SOURCES: docs/claude_project/05_RESULTS.md  (every measured number)
%          docs/claude_project/08_OPEN_ITEMS.md §4 and §6 (the gaps and risks)
%          build/impl_*/ and build/bd_187p5_explore/ via numbers.tex
%
% GAPS before this chapter can be written honestly:
%   1. ACCURACY is measured (DD-034): float and int8 on the held-out test split
%      of data/grape_leafsplit, split by leaf, scored once; every figure is a
%      \thAcc* macro from software/runs/mobilenetv2_grape_leafsplit_test.json.
%      STILL OPEN: it was measured on the retrained checkpoint \thAccCkptSha,
%      while software/export/ - and so every bit-exactness result - is built
%      from \thCkptSha. Until the export, the golden vectors and the
%      simulations are regenerated from \thAccCkptSha, the two must not be
%      presented as one model; numbers.tex carries a WARNING line while they
%      differ.
%   2. The \MHz{\thFmaxPhaseFour} figure is unattributed (Section~\ref{sec:phase4}).
% ===========================================================================
\chapter{Experimental Evaluation}
\label{chap:evaluation}

\section{Overall Performance and Resource Utilization}
\label{sec:overall}
% The headline table. Separate OUT-OF-CONTEXT from IN-CONTEXT explicitly -
% they measure different things and mixing them would misreport your own data:
%
%   out of context, after the image path:
%     accel_top  \MHz{\thFmaxAccelTop}  \thLutAccelTop{} LUT
%     accel_sys  \MHz{\thFmaxAccelSys}  \thLutAccelSys{} LUT
%     accel_axi  \MHz{\thFmaxAccelAxi}  \thLutAccelAxi{} LUT
%   in context, ZCU104 block design:
%     \MHz{\thClkBoard} constrained, WNS +\thWnsBoard{} ns, \thLutBoard{} LUT
%
% DSP is \thDspAccelAxi{} (\thDspPctAccelAxi\%) throughout; URAM \thUramAccelAxi{}.

\subsection{Fmax Is a Band, Not a Point}
\label{sec:fmax-band}
% One of the more honest results in the thesis, so give it its own subsection.
% Across seven post-route runs the figure spans roughly 197-208 MHz. The
% ordering is NON-MONOTONIC IN LOGIC SIZE - the largest design is the fastest -
% so the spread is the placer, not a cost of the DMA and the CSR. accel_top
% reproduced EXACTLY on a re-run, which shows the flow is deterministic and the
% spread is this design's sensitivity to small RTL changes rather than tool
% noise: accel_sys moved +9.3 MHz from a change that only turned a port into an
% internal net.
% Conclusion to state: no point estimate from this flow is defensible, and all
% seven launch from the same act_buffer URAM read (Section~\ref{sec:dd023}) - only the
% endpoint moves.

\section{Sequential Optimization Impact Analysis}
\label{sec:opt-impact}
% Chapter~\ref{chap:timing} presents each change with its measurement; this section is
% the aggregate view - the cumulative curve and the cost in LUT/FF per MHz
% gained. Do not repeat the narrative, reference it.

\section{Throughput and Latency}
\label{sec:throughput}
% \thCycles{} cycles per inference, measured (feeder-busy) in simulation.
% At \MHz{\thClkBoard}: \thLatencyBoard{} ms. Give the cycle model's prediction
% alongside and report the error - it came within 7%, which is worth saying
% because it is a validation of the model in Section~\ref{sec:roofline-analysis}.
% Weight loading is \thDmaLoadCycles{} cycles against \thCyclesCompute{} of
% compute: \thLoadOverheadPct\% overhead, none of it hidden (C-11).

\section{Correctness and Accuracy Assessment}
\label{sec:correctness}

\subsection{Hardware--Software Bit-Exactness}
\label{sec:bit-exactness-result}
% \thElements elements over \thInstructions instructions, exact, on
% \thImagesVerified{} images - one per class, all four argmax outcomes
% exercised, all ALL PASS. \thTestbenches/\thTestbenches{} testbenches passing.
% This is the thesis's central result. Chapter~\ref{chap:verification} is the method;
% this is the number.

\subsection{Classification Accuracy}
\label{sec:accuracy}
% The protocol first, in two sentences: trained on train, checkpoint chosen on
% val (epoch \thAccBestEpoch), int8 scales calibrated on train, and only then
% the test split scored - once. Then the result, from numbers.tex:
%
%   float \thAccFloat\%, int8 \thAccInt\% on \thAccTestImages{} test images of
%   \thAccTestLeaves{} leaves; drop \thAccDrop{} points; int8 and float
%   disagree on \thAccDisagree{} images.
%
% Quote BOTH 95% intervals and say why there are two: the photos of a leaf
% are not independent, so "per image" [\thAccIntCiImageLo, \thAccIntCiImageHi]
% is too narrow and "per leaf" [\thAccIntCiLeafLo, \thAccIntCiLeafHi] is
% conservative. Per-class precision / recall / F1 and macro-F1
% (\thAccIntMacroF) as a table, and the confusion matrix (\thAccIntConf*).
%
% The int8 figure is the INTEGER PIPELINE's, not fake-quant's: the arithmetic
% the RTL reproduces bit for bit (Section~\ref{sec:bit-exactness-result}), checked
% against the accumulator bound of Section~\ref{sec:acc-bound}. So it is the
% accelerator's accuracy - for the checkpoint it was measured on (see the
% header of this file).
%
% Report the retracted number too, as the thing that was wrong and how it was
% found (Section~\ref{sec:leakage}): 100.00% on data/grape (rotated copies) -
% and that the first fix, grape_clean, leaked at the level of the leaf, so
% nothing measured on it is quoted either. A thesis that explains a retracted
% number is stronger than one that never had it.
% Separately: export.py --check 32 gives 32/32 agreement between the integer
% pipeline and both the float and fake-quant models. That measures the
% QUANTISATION, not the model, and the distinction should be explicit.

\section{Power Consumption}
\label{sec:power}
% Vivado's post-route estimate, clearly labelled ESTIMATED (Section~\ref{sec:overall}
% is measured; do not let the two sit in one table without a marker).
% Energy per inference is the comparable quantity for Section~\ref{sec:comparison}.

\section{Comparison with Baselines and Related Work}
\label{sec:comparison}
% RENAMED from "Comparison with the Reference Implementation", which was
% ambiguous and too narrow. Three comparisons, in increasing order of risk:
%   1. the bit-exact C model on the host CPU (Section~\ref{sec:c-model}) - a fair,
%      like-for-like baseline because it is the SAME arithmetic;
%   2. the vendor DPU [xilinx_pg338] on the same device class - the "why not
%      Vitis-AI?" question, which will be asked;
%   3. the FPGA accelerators of Section~\ref{sec:related-accel} - and here state the
%      caveat before the table, not after: different networks, devices,
%      precisions and reporting conventions make GOP/s comparisons weak. Say
%      what is comparable (DSP efficiency, MAC/cycle/DSP) and what is not.
